\chapter*{前言与版本边界}
\addcontentsline{toc}{chapter}{前言与版本边界}
本工程面向掌握多元微积分、线性代数和初步点集拓扑的读者，以流形微积分和丛联络为主线编写中文教材初稿。主参考为周建伟《微分几何讲义》；Chern、Chen、Lam的英文讲义供交叉核对，梁灿彬、周彬的著作供几何直觉与相对论参考。

第1--2章写入拓扑、欧氏分析、对偶与外代数预备，以及曲线曲面和Gauss内蕴曲率；第1章配梁书16题改述及6道补充题，第2章配8题。第3章建立光滑流形基础并配6题，第4章建立切空间、切映射及切丛/余切丛。原题、改编和原创补充分开登记。

第5--9章依次建立子流形与局部流、张量与向量丛、微分形式、积分与Stokes、de Rham上同调。第10--13章讨论一般联络、黎曼几何、完备性与变分、子流形相容方程和极小图像。第14章补入Sobolev与弱解工具，证明Hodge分解及曲面Gauss--Bonnet的有限剖分版本。

第15章建立李群与齐性几何，第16章构造主丛、相伴丛及规范变换，第17章讨论复图册、Hermite/Kahler、Chern联络与Ricci形式。第18章建立Chern--Weil闭形式、联络无关性和相对传递，构造实Chern/Pontryagin/Euler类并计算射影示性数，区分局部Chern--Simons及怪球专题的外引范围。第19章构造Clifford代数与Spin双覆盖，建立旋量联络、Dirac平方及标量曲率消失结论，并证明Kahler/平坦特殊Lagrange校准的具体应用。第4--19章各配8题详解，来源逐题登记。

第20章补证紧致无边界情形的Dirac型算子Fredholm性质与指标稳定性，推导Yang--Mills变分、四维能量界及ASD规范固定的椭圆符号，并在给定Spin$^c$数据下写出Seiberg--Witten方程和非负标量曲率消失结论；另配8题详解。Atiyah--Singer公式仅精确外引，其完整证明、Spin$^c$拓扑比较、模空间紧性与不变量构造仍未建立。第21--23章为下一写作计划。

已有正文也保留明确边界：分析以实数完备性、单变量微积分、有限维矩阵理论为先修；Hodge部分另用Lebesgue积分和L²完备性，第20章一般Hilbert指标稳定性另将有界逆定理作为分析先修。曲面剖分存在/拓扑比较、一般Chern--Gauss--Bonnet整体证明、NN充分性、一般局部复势、整示性类比较、Spin障碍与拓扑$w_2$比较及怪球构造分类均未在本书建立。球面复结构的2026年外部研究只登记核查范围，不充作本书已证定理。

本版含63幅TikZ图。来源补核与独立数学审阅尚未完成；图示、源码、模型及编译检查均不代替数学审阅。
